% Part: first-order-logic
% Chapter: tableaux
% Section: quantifier-rules

\documentclass[../../../include/open-logic-section]{subfiles}

\begin{document}

\olfileid{fol}{tab}{qrl}

\olsection{د سورونو قاعدې}

\subsection{د~$\lforall$ قاعدې}

\begin{defish}
\AxiomC{\sFmla{\True}{\lforall[x][!A(x)]}}
\RightLabel{\TRule{\True}{\forall}}
\UnaryInfC{\sFmla{\True}{!A(t)}}
\DisplayProof
\hfill
\AxiomC{\sFmla{\False}{\lforall[x][!A(x)]}}
\RightLabel{\TRule{\False}{\lforall}}
\UnaryInfC{\sFmla{\False}{!A(a)}}
\DisplayProof
\end{defish}

په~\TRule{\True}{\lforall} کښې $t$ تړلې اصطلاح ده (يعنې متغيرونه
نۀ لري)۔ په~\TRule{\False}{\lforall} کښې $a$ داسې !!a{constant}
دے چې بايد د~\TRule{\False}{\lforall} تر قاعدې پاس په څانګه کښې
هېڅ ځاے نۀ وي راغلے۔ $a$ د~\TRule{\False}{\forall} استنتاج
\emph{ځانګړے متغير} بولو۔\footnote{د «ځانګړي متغير» اصطلاح کاروو،
سره له دې چې په پورته قاعده کښې $a$ يو !!a{constant} دے۔ دا
تاريخي لاملونه لري۔}

\subsection{د~$\lexists$ قاعدې}

\begin{defish}
\AxiomC{\sFmla{\True}{\lexists[x][!A(x)]}}
\RightLabel{\TRule{\True}{\lexists}}
\UnaryInfC{\sFmla{\True}{!A(a)}}
\DisplayProof
\hfill
\AxiomC{\sFmla{\False}{\lexists[x][!A(x)]}}
\RightLabel{\TRule{\False}{\lexists}}
\UnaryInfC{\sFmla{\False}{!A(t)}}
\DisplayProof
\end{defish}

بيا هم $t$ تړلې اصطلاح ده، او~$a$ داسې !!a{constant} دے چې د
\TRule{\True}{\lexists} تر قاعدې پاس په څانګه کښې نۀ راځي۔
$a$ د~\TRule{\True}{\lexists} استنتاج \emph{ځانګړے متغير} بولو۔

هغه شرط چې ځانګړے متغير د~\TRule{\False}{\lforall} يا
\TRule{\True}{\lexists} استنتاج تر پاسه په څانګه کښې نۀ راځي،
\emph{د ځانګړي متغير شرط} بلل کېږي۔

\begin{explain}
هغه قرارداد را ياد کړئ چې کله $!A$ داسې !!a{formula} وي چې
!!{variable}~$x$ پکې ازاد وي، دا د~$!A(x)$ په ليکلو ښيو۔ په همدې
سياق کښې $!A(t)$ د~$\Subst{!A}{t}{x}$ لنډه بڼه ده۔ نو د
\TRule{\False}{\lexists} قاعده داسې هم ليکلے شو:
\begin{prooftree}
  \AxiomC{\sFmla{\False}{\lexists[x][!A]}}
  \RightLabel{\TRule{\False}{\lexists}}
  \UnaryInfC{\sFmla{\False}{\Subst{!A}{t}{x}}}
\end{prooftree}
پام وکړئ چې $t$ له مخکښې په~$!A$ کښې راتلے شي؛ مثلاً $!A$ کېدے
شي $\Atom{\Obj P}{t,x}$ وي۔ نو له
$ \sFmla{\False}{\lexists[x][\Atom{\Obj P}{t,x}]}$ څخه
$\sFmla{\False}{\Atom{\Obj P}{t,t}}$ استنتاجول د
\TRule{\False}{\lexists} سمه کارونه ده۔ خو په
\TRule{\False}{\lforall} او~\TRule{\True}{\lexists} کښې د ځانګړي
متغير شرطونه غواړي چې !!{constant}~$a$ په~$!A$ کښې نۀ راځي۔ نو له
$\sFmla{\False}{\lforall[x][\Atom{\Obj P}{a,x}]}$
څخه د~$\TRule{\False}{\lforall}$ په کارولو
$\sFmla{\False}{\Atom{\Obj P}{a,a}}$ په سمه توګه نۀ شئ استنتاجولے۔
\end{explain}

\begin{explain}
په~\TRule{\True}{\lforall} او~\TRule{\False}{\lexists} کښې د
تړلې اصطلاح~$t$ پر انتخاب نور بنديز نشته۔ بل لوري ته، د
\TRule{\True}{\lexists} او~\TRule{\False}{\lforall} په قاعدو کښې
د ځانګړي متغير شرط غواړي چې !!{constant}~$a$ د اړوند استنتاج تر
پاسه په څانګو کښې هېڅ ځاے نۀ راځي۔ دا شرط د نظام د سموالي د
يقيني کولو دپاره اړين دے۔ له دې شرط پرته لاندې به د
$\lexists[x][\formula{A}(x)] \lif \lforall[x][\formula{A}(x)]$
دپاره يو تړلے !!{tableau} و:
\begin{center}
\begin{tableau}{}
  [\sFmla{\False}{\lexists[x][\formula{A}(x)] \lif \lforall[x][\formula{A}(x)]}, just=\TAss
    [\sFmla{\True}{\lexists[x][\formula{A}(x)]},
      just={\TRule{\False}{\lif}[1]}
      [\sFmla{\False}{\lforall[x][\formula{A}(x)]},
        just={\TRule{\False}{\lif}[1]}
        [\sFmla{\True}{\formula{A}(a)},
          just={\TRule{\True}{\lexists}[2]}
          [\sFmla{\False}{\formula{A}(a)},
            just={\TRule{\False}{\lforall}[3]}, close]
        ]
      ]
    ]
  ]
\end{tableau}
\end{center}
خو $\lexists[x][\formula{A}(x)] \lif
\lforall[x][\formula{A}(x)]$ معتبر نۀ دے۔
\end{explain}

\end{document}
